\documentclass{article}
\usepackage{amsmath,amssymb}
\usepackage{xurl}
\usepackage[hidelinks]{hyperref}
\setlength{\emergencystretch}{2em}
% Shared commands for notes in docs/paper-gaps/.

\DeclareUnicodeCharacter{2080}{\ifmmode{}_{0}\else\textsubscript{0}\fi}
\DeclareUnicodeCharacter{2081}{\ifmmode{}_{1}\else\textsubscript{1}\fi}
\DeclareUnicodeCharacter{2082}{\ifmmode{}_{2}\else\textsubscript{2}\fi}
\DeclareUnicodeCharacter{2099}{\ifmmode{}_{n}\else\textsubscript{n}\fi}

\newcommand{\Lean}{\textsc{Lean}}
\ExplSyntaxOn
\NewDocumentCommand{\leanid}{m}{
  \ifmmode
    \textsf{\detokenize{#1}}
  \else
    \group_begin:
    \urlstyle{sf}
    \tl_set:Nn \l_tmpa_tl {#1}
    \tl_replace_all:Nnn \l_tmpa_tl {\_} {_}
    \exp_args:NV \path \l_tmpa_tl
    \group_end:
  \fi
}
\ExplSyntaxOff
\providecommand{\tr}{\operatorname{tr}}
\newcommand{\tnleanIssueBase}{https://github.com/LionSR/TNLean/issues/}
\newcommand{\tnleanPrBase}{https://github.com/LionSR/TNLean/pull/}
\newcommand{\tnleanDocBase}{https://sirui-lu.com/TNLean/docs/}
\newcommand{\ghissue}[1]{\href{\tnleanIssueBase#1}{GitHub issue \##1}}
\newcommand{\ghpr}[1]{\href{\tnleanPrBase#1}{PR \##1}}
\newcommand{\doclink}[1]{\href{\tnleanDocBase#1.html}{\path{#1}}}

% Dirac notation and the identity operator, as in blueprint/src/macros/common.tex.
% \providecommand keeps note-local definitions authoritative.
\usepackage{dsfont}
\providecommand{\ket}[1]{|#1\rangle}
\providecommand{\bra}[1]{\langle #1|}
\providecommand{\braket}[2]{\langle #1 | #2 \rangle}
\providecommand{\ketbra}[2]{|#1\rangle\!\langle #2|}
\providecommand{\Id}{\mathds{1}}
\providecommand{\one}{\mathds{1}}
\providecommand{\C}{\mathbb{C}}
\providecommand{\MN}[1]{M_{#1}(\C)}
\providecommand{\MD}{M_D(\C)}

% Verdict marker: \gapnote{<kind>}{<status>}. Kind is the policy
% classification (clarification, local-correction, scope-restriction,
% unfaithful, false-source, open-gap); status is open, wip (elimination
% underway), resolved, or historical. Machine-read by the site index;
% prints a verdict line.
\newcommand{\gapnote}[2]{\par\noindent\textbf{Verdict:} #1 (#2).\par}

\title{The Mixed Interpolation Path: Commutation of the Canonical Parent with the Symmetry}
\author{}
\date{2026-11-16}
\begin{document}
\maketitle
\gapnote{scope-restriction}{open}

\section{The source claim}
Garre-Rubio, Lootens and Moln\'ar connect two injective matrix product
states with equal $L$-symbols by an interpolating tensor $A(\gamma)$ and
common operators $T_a$ \cite[Section~5]{GarreRubio2023MPOSym}. The
section opens with
\begin{quote}
    ``We do so by explicitly constructing a continuous and symmetric
    embedding path that connects both MPSs and therefore, their parent
    Hamiltonians.''
\end{quote}
and the path Hamiltonian is fixed as the canonical parent:
\begin{quote}
    ``It is clear that the tensor $A(\gamma)$ defined in
    Eq.~(defAgamma) is injective for $\gamma \in (0,1)$, so that the
    parent Hamiltonian $H_{A(\gamma)} \equiv H(\gamma)$ is gapped and has
    as unique ground state $|\psi_{A(\gamma)} \rangle$ within that
    region. [\ldots] We remind that the local Hamiltonian terms are
    defined by $h_\gamma = \mathbb 1 - \Pi_{S_\gamma}$.''
\end{quote}
The definition of a phase requires that ``the Hamiltonian path commutes
with an injective MPO representation of $\mathcal C$'', and the
definition of symmetry states that ``the unitarity of the operators is
not required''.\footnote{Source traceability:
    \path{Papers/2203.12563/REsubmission.tex}, line 1582 (symmetric
    embedding path), lines 1687--1690 (the canonical parent $H(\gamma)$
    and $h_\gamma=\mathbb 1-\Pi_{S_\gamma}$), lines 1240--1246
    (definition \texttt{defsym} and the remark on unitarity) and
    lines 1253--1259 (definition \texttt{defphase}).}
Section~5 proves exact covariance of $A(\gamma)$ under $T_a$
(\texttt{Agammasym}) and the fusion of the $T_a$, but never shows that
$H(\gamma)$ commutes with the periodic operators $O_a$.

The symmetrization of the parent Hamiltonian, lines 1290--1320, is the
source's tool for commutation: the local term is averaged with the
canonical left integral $\Omega$ of a $C^*$-weak Hopf algebra and its
pulling-through tensors. Section~5 does not invoke it for the path, and
it does not construct an integral, a coproduct, or a star structure for
the common operators $T_a$; it constructs only their fusion tensors and
$F$- and $L$-symbols.

\section{The canonical parent need not commute}
Let $S\subset(\mathbb C^d)^{\otimes2}$ be the two-site support of an
injective tensor $A$, $P$ its orthogonal projection, and $h=\mathbb 1-P$.
Covariance under a symmetry $O$ gives $O\,S\subset S$. Commutation
$[h,O]=0$ is equivalent to invariance of both $S$ and $S^\perp$, and the
latter is equivalent to $O^\dagger S\subset S$. For unitary or
adjoint-closed symmetries this is automatic; covariance alone does not
supply it.

\paragraph{Explicit data.} Take $d=4$, bond dimension $2$, and the tensor
\[
    A^1=E_{11},\qquad A^2=E_{12},\qquad A^3=E_{21},\qquad
    A^4=\mathbb 1_2,
\]
which is one-site injective since its letters span $\mathcal M_2$. Put
\[
    V=\begin{pmatrix}1&2\\0&-1\end{pmatrix},\qquad V^2=\mathbb 1,
    \qquad
    u=\begin{pmatrix}
        1 & 2  & 0  & 0  \\
        0 & -1 & 0  & 0  \\
        4 & 4  & -1 & -2 \\
        0 & 0  & 0  & 1
    \end{pmatrix}.
\]
Then $\sum_l u_{kl}A^l=VA^kV^{-1}$ for every $k$; for instance
$VE_{11}V=E_{11}+2E_{12}$ is the first row. Also $u^2=\mathbb 1$, but $u$ is
not unitary.

\paragraph{The Section~5 data are present.} Take $\mathcal C=\mathbb Z_2$
with $O_e=\mathbb 1$ and $O_g=u^{\otimes N}$, both of bond dimension one.
Fusion is exact, since $u^2=\mathbb 1$. The action of $O_g$ on $A$ is exact,
with one copy and the action maps $W_g=V$ and $V_g=V^{-1}=V$, which satisfy
$V_gW_g=\mathbb 1$. Choose $A_0=A_1=A$ and identical operator data at both
endpoints. Then the chosen raw $L$ matrices, and hence the $L$- and
$F$-symbols, agree. These are exactly the hypotheses under which the
common operator of the path is formalized.

\paragraph{Failure of $S^\perp$-invariance.} The support is
$S=\{s_X : X\in\mathcal M_2\}$ with
$(s_X)_{kl}=\operatorname{tr}(A^kA^lX)$, and
$(u\otimes u)s_X=s_{VXV}$, so $S$ is invariant. Let $w=e_2\otimes e_1$.
Since $A^2A^1=E_{12}E_{11}=0$, every $(s_X)_{21}=0$, so $w\in S^\perp$.
Put $a=ue_2=(2,-1,4,0)$ and $b=ue_1=(1,0,4,0)$, the second and first
columns of $u$, so that $(u\otimes u)w=a\otimes b$. All data are real, so
\[
    \langle s_{E_{11}},(u\otimes u)w\rangle
    =\Bigl(\bigl(\textstyle\sum_ka_kA^k\bigr)
    \bigl(\textstyle\sum_lb_lA^l\bigr)\Bigr)_{11}
    =\left[\begin{pmatrix}2&-1\\4&0\end{pmatrix}
        \begin{pmatrix}1&0\\4&0\end{pmatrix}\right]_{11}
    =-2\neq0.
\]
Hence $(u\otimes u)S^\perp\not\subset S^\perp$, and
$[h,u\otimes u]\neq0$. The Gram matrix of the basis
$s_{E_{11}},s_{E_{12}},s_{E_{21}},s_{E_{22}}$ is
\[
    \begin{pmatrix}5&0&0&1\\0&3&0&0\\0&0&3&0\\1&0&0&2\end{pmatrix},
\]
so $P$ and $h$ have rational entries. The periodic Hamiltonian
$H=\sum_ih_i$ on three sites has the rational commutator entry
$\langle 111|[H,u^{\otimes3}]|112\rangle=2/9$, so the global
commutation fails as well.

\paragraph{Consequence for the path.} By the definition of a phase,
$H(0)$ restricts to the endpoint Hamiltonian $H_0$ on
$\mathcal H_0^{\otimes N}$, and $O_a=O^0_a\oplus O^1_a\oplus O^{\rm path}_a$.
If the canonical $H(\gamma)$ commuted with $O_a$ at $\gamma=0$, the
restriction would give $[H_0,O^0_g]=0$, where $H_0$ is the canonical
parent of $A$ and $O^0_g=u^{\otimes N}$. This was refuted above. Thus the
source's claim that the canonical parents form a symmetric path does not
follow from the data that Section~5 constructs.

\section{The formal boundary}
The enlarged mixed interaction is the canonical two-site parent of the
enlarged mixed tensor on the open interval. It is continuous through both
endpoints. Suppose the two-site arbitrary-boundary range
$\{O^2_{\widetilde T,X}:X\in\mathcal M_\chi\}$ of the common operator is
closed under the adjoint. Then the interaction commutes with every
two-site boundary operator for every $\gamma\in[0,1]$. Hence the periodic
path Hamiltonian commutes with the periodic common operator and with each
boundary operator whose boundary commutes with its letters, at every ring
length $N\geq2$.\footnote{\leanid{MPSTensor.MPOSymmetry.arbitraryPhysicalMixedInteraction_commute_mpoWithBoundary}
    and
    \leanid{MPSTensor.MPOSymmetry.interactionHamiltonian_arbitraryPhysicalMixedInteraction_commute_mpo},
    in \doclink{TNLean/MPS/Symmetry/MPOSymmetry/ArbitraryPhysicalMixedParentCommutation}.}
The fusion hypotheses and the equal $L$ matrices are not used for this
conclusion. Adjoint closure is the only hypothesis added to those of the
action decompositions and one-site injectivity, and the example above
shows that it cannot be dropped. The source-labelled claim of a symmetric
canonical path remains unformalized.

\section{Elimination plan}
\begin{enumerate}
    \item The adjoint-closed commutation above is formalized and is the
          scope-restricted replacement for the canonical path.
    \item Replace the canonical parent by the symmetrized parent of
          lines 1290--1320 once the weak-Hopf averaging route is
          formalized and the common operators are shown to represent a
          $C^*$-weak Hopf algebra; see the separate review of support
          averaging versus ambient completion,
          \texttt{glm23\_wha\_parent\_completion.tex}.
    \item Derive adjoint closure of the common operator from physical
          adjoint data of the endpoint families, in the manner of the
          normal periodic adjoint duality already used for single
          compatible tensors. For unitary on-site group symmetries this
          closure holds automatically at each endpoint; whether it
          survives the mixed off-diagonal blocks is not yet formalized.
\end{enumerate}

\bibliographystyle{alpha}
\bibliography{references}
\end{document}
